% =============================================================================
% 06_fenicsx_cfd_implementation.tex
%
% FEniCSx / dolfinx 0.10 Implementation Guide
% Stokes CFD on AnXplore geometries — WSS ground-truth pipeline
%
% Purpose: authoritative implementation reference for Claude Code.
% All API patterns verified against dolfinx v0.10.0.post0 official docs.
% Do NOT use patterns from older dolfinx versions (0.4--0.8 have
% incompatible import paths and function signatures).
%
% Read this document when: writing or reviewing any code in
%   src/hemodyn_pinn/cfd/stokes_solver.py
%   src/hemodyn_pinn/cfd/ns_solver.py
%   src/hemodyn_pinn/cfd/postprocess.py
%   scripts/02_run_cfd.py
% =============================================================================

\documentclass[11pt]{article}
\usepackage[a4paper, margin=1in]{geometry}
\usepackage{amsmath, amssymb}
\usepackage{hyperref}
\usepackage{xcolor}
\usepackage{enumitem}
\usepackage{titlesec}
\usepackage{listings}
\usepackage{booktabs}
\hypersetup{colorlinks=true, linkcolor=blue!50!black, urlcolor=blue!50!black}
\titleformat{\section}{\large\bfseries\color{blue!40!black}}{\thesection.}{0.5em}{}
\titleformat{\subsection}{\normalsize\bfseries}{\thesubsection}{0.5em}{}

\definecolor{codebg}{RGB}{245,245,242}
\lstset{
  language=Python,
  basicstyle=\ttfamily\small,
  backgroundcolor=\color{codebg},
  frame=single, framerule=0pt, rulecolor=\color{codebg},
  breaklines=true, breakatwhitespace=true,
  showstringspaces=false,
  commentstyle=\color{gray},
  keywordstyle=\color{blue!70!black},
}

\title{\bfseries FEniCSx CFD Implementation Guide\\[0.3em]
\large Stokes solver, boundary tagging, WSS extraction, and output format\\
for the WSS-from-MRI PINN project --- dolfinx 0.10}
\author{}
\date{}

\begin{document}
\maketitle

\section*{How to use this document}

This is a precision engineering reference, not a tutorial. It records
the exact dolfinx~0.10 API patterns required for the CFD ground-truth
pipeline, the gotchas specific to this project's geometry, and the
decisions already made. Each section states what to do, gives the
verified code pattern, and explains why (because silent API mistakes in
FEniCSx produce solutions that look correct but are wrong).

\textbf{Version lock.} All patterns are verified against
\texttt{dolfinx v0.10.0.post0}. Do not use patterns from older versions;
the import paths (\texttt{basix.ufl}, \texttt{dolfinx.fem.petsc}),
function names (\texttt{functionspace} vs \texttt{FunctionSpace}), and
assembly API changed substantially between 0.7 and 0.9.

\section{The strong form and its physical parameters}

The steady Stokes equations in physical (dimensional) variables:
\begin{align}
-\mu\,\Delta\mathbf{u} + \nabla p &= \mathbf{0} \quad \text{in } \Omega, \label{eq:stokes_mom}\\
\nabla\cdot\mathbf{u} &= 0 \quad \text{in } \Omega. \label{eq:stokes_cont}
\end{align}

\textbf{Note on sign convention.} The dolfinx demo uses the symmetrised form
$-\nabla\cdot(\nabla\mathbf{u} + p\mathbf{I}) = \mathbf{f}$, which flips the
sign of the pressure relative to the classical definition. This gives a
symmetric saddle-point system. The sign flip means the pressure solution
from dolfinx satisfies $p_{\text{dolfinx}} = -p_{\text{physical}}$; correct
for this when reporting or exporting pressure.

\textbf{Physical parameters (Phase~1--2, Newtonian blood):}
\begin{itemize}[noitemsep]
  \item Density: $\rho = 1060\,\text{kg/m}^3$ (not needed in Stokes, only for Re check)
  \item Dynamic viscosity: $\mu = 3.5\times10^{-3}\,\text{Pa}\cdot\text{s}$
  \item Inlet mean velocity: $U = 2\times10^{-5}\,\text{m/s}$
    (gives $\mathrm{Re} = \rho UL/\mu \approx 0.10$ with $L = 0.01628\,\text{m}$;
    confirmed Stokes regime)
  \item Mesh coordinates: millimetres. Convert to metres before setting up any
    FEniCSx form: multiply all coordinate values by \texttt{1e-3}.
\end{itemize}

\section{Mesh loading: from VTK to dolfinx}

AnXplore meshes are tetrahedral VTK files (\texttt{Fluid\_*.vtk}).
dolfinx does not read VTK directly; convert via meshio to XDMF+H5 first.

\begin{lstlisting}
import meshio
import numpy as np

def vtk_to_xdmf(vtk_path: str, xdmf_path: str) -> None:
    """Convert an AnXplore VTK mesh to XDMF+H5 for dolfinx.

    Scales coordinates from mm to m and writes a single-domain
    tetrahedral mesh. Surface triangles (if present) are dropped;
    boundary facets are detected topologically in dolfinx.
    """
    mesh = meshio.read(vtk_path)
    # Scale mm -> m
    mesh.points = mesh.points * 1e-3
    # Keep only tetrahedral cells; drop triangles and anything else
    tet_cells = [c for c in mesh.cells if c.type == "tetra"]
    out = meshio.Mesh(points=mesh.points, cells=tet_cells)
    meshio.write(xdmf_path, out)
\end{lstlisting}

\textbf{Why drop surface triangles?} dolfinx builds its own boundary
facet topology from the tet connectivity. Importing mixed cell-type meshes
(tetra + triangle) confuses dolfinx's topology layer and causes incorrect
boundary detection. Extract boundary facets inside dolfinx instead
(see Section~\ref{sec:boundary_tagging}).

\subsection*{Loading the XDMF mesh in dolfinx}

\begin{lstlisting}
from mpi4py import MPI
from dolfinx.io import XDMFFile
from dolfinx.mesh import Mesh

def load_mesh(xdmf_path: str) -> Mesh:
    with XDMFFile(MPI.COMM_WORLD, xdmf_path, "r") as f:
        mesh = f.read_mesh(name="Grid")
    # Always create connectivity for facet-level operations
    mesh.topology.create_connectivity(
        mesh.topology.dim - 1, mesh.topology.dim
    )
    return mesh
\end{lstlisting}

The \texttt{create\_connectivity} call is mandatory before any boundary
tagging. Omitting it causes silent empty facet sets and hard-to-diagnose
boundary condition failures.

\section{Boundary tagging}
\label{sec:boundary_tagging}

AnXplore geometries have three boundary types that must be tagged
with integer markers before imposing BCs:

\begin{center}
\begin{tabular}{lll}
\toprule
\textbf{Tag} & \textbf{Boundary} & \textbf{Detection criterion} \\
\midrule
1 & Wall $\Gamma_w$ & All boundary facets not tagged as inlet or outlet \\
2 & Inlet $\Gamma_\text{in}$ & Facets at $y \approx 0$, centroid $x < 0$ \\
3 & Outlet $\Gamma_\text{out}$ & Facets at $y \approx 0$, centroid $x > 0$ \\
\bottomrule
\end{tabular}
\end{center}

This topology was confirmed by mesh inspection (CLAUDE.md open question~4,
now closed): both inlet and outlet are circular disks at $y=0$ with outward
normal $(0,-1,0)$; they are separated at $x=0$. Tube radius $\approx 2\,\text{mm}$
$= 0.002\,\text{m}$ after scaling.

\begin{lstlisting}
import numpy as np
from dolfinx.mesh import locate_entities_boundary, meshtags

def tag_boundaries(mesh) -> "dolfinx.mesh.MeshTags":
    """Return a MeshTags object with integer tags on all boundary facets.

    Tags: 1=wall, 2=inlet (x<0, y~0), 3=outlet (x>0, y~0).
    All coordinates are in metres (mesh already scaled mm->m).
    """
    fdim = mesh.topology.dim - 1
    tol = 1e-6  # geometric tolerance in metres

    def inlet(x):
        return np.isclose(x[1], 0.0, atol=tol) & (x[0] < 0.0)

    def outlet(x):
        return np.isclose(x[1], 0.0, atol=tol) & (x[0] >= 0.0)

    # Locate facets for inlet and outlet
    inlet_facets  = locate_entities_boundary(mesh, fdim, inlet)
    outlet_facets = locate_entities_boundary(mesh, fdim, outlet)

    # Collect all boundary facets, then assign wall tag to the remainder
    all_boundary = locate_entities_boundary(mesh, fdim, lambda x: np.ones(x.shape[1], dtype=bool))

    inlet_set  = set(inlet_facets.tolist())
    outlet_set = set(outlet_facets.tolist())
    wall_facets = np.array(
        [f for f in all_boundary if f not in inlet_set and f not in outlet_set],
        dtype=np.int32
    )

    # Build combined sorted index/values arrays for meshtags
    indices = np.concatenate([wall_facets, inlet_facets, outlet_facets])
    values  = np.concatenate([
        np.full(len(wall_facets),   1, dtype=np.int32),
        np.full(len(inlet_facets),  2, dtype=np.int32),
        np.full(len(outlet_facets), 3, dtype=np.int32),
    ])
    sort_idx = np.argsort(indices)
    return meshtags(mesh, fdim, indices[sort_idx], values[sort_idx])
\end{lstlisting}

\textbf{Sanity check (always run):} after tagging, assert
\texttt{np.sum(tags.values == 2) > 0} and
\texttt{np.sum(tags.values == 3) > 0}. Empty inlet or outlet tags mean
no-slip will be applied everywhere and the solve will return
$\mathbf{u} \equiv \mathbf{0}$.

\section{Function spaces: Taylor--Hood elements}

Taylor--Hood ($\mathbf{P}_2 / P_1$) is the standard stable pair for Stokes:
piecewise quadratic velocity, piecewise linear pressure. It satisfies the
inf-sup (LBB) condition on tetrahedral meshes.

\begin{lstlisting}
import ufl
from basix.ufl import element, mixed_element
from dolfinx import default_real_type
from dolfinx.fem import functionspace, Function

def make_function_spaces(mesh):
    """Create Taylor-Hood P2/P1 velocity and pressure spaces.

    Returns (V, Q) as separate spaces for the blocked formulation,
    which is preferred for iterative solvers.
    """
    P2 = element(
        "Lagrange", mesh.basix_cell(), degree=2,
        shape=(mesh.geometry.dim,), dtype=default_real_type
    )
    P1 = element(
        "Lagrange", mesh.basix_cell(), degree=1,
        dtype=default_real_type
    )
    V = functionspace(mesh, P2)  # velocity: vector P2
    Q = functionspace(mesh, P1)  # pressure: scalar P1
    return V, Q
\end{lstlisting}

\textbf{Why separate spaces, not mixed?} The blocked ($V \times Q$
separately) formulation is preferred here because it allows the MINRES +
block-diagonal preconditioner (see Section~\ref{sec:solver}), which
converges much faster than a direct solver on meshes of ~160k nodes.
The mixed element approach (\texttt{mixed\_element([P2, P1])}) is
simpler but slower for large problems; use it only for debugging.

\section{Boundary conditions}

\subsection*{Inlet: parabolic Poiseuille profile}

For a circular inlet of radius $R = 0.002\,\text{m}$ centred at the
centroid of the inlet facets, the Poiseuille profile is:
\[
  u_y(r) = U_{\max}\!\left(1 - \frac{r^2}{R^2}\right), \qquad
  U_{\max} = 2\bar{U}, \qquad
  r = \sqrt{(x - x_c)^2 + (z - z_c)^2}.
\]
Velocity direction: $(0, -1, 0)$ (into the domain, outward normal is
$(0, -1, 0)$ pointing out, so flow enters in the $-y$ direction).

\begin{lstlisting}
import numpy as np
from dolfinx.fem import Function, locate_dofs_topological, dirichletbc

def make_inlet_bc(V, mesh, tags, U_mean: float = 2e-5):
    """Parabolic Poiseuille inlet BC on the inlet facets (tag=2).

    Parameters
    ----------
    U_mean : float
        Mean inlet velocity in m/s. Default 2e-5 gives Re ~ 0.10.
    """
    R = 0.002        # tube radius in metres
    U_max = 2.0 * U_mean

    # Find centroid of inlet facets to locate tube axis
    fdim = mesh.topology.dim - 1
    inlet_facets = tags.find(2)
    inlet_dofs = locate_dofs_topological(V, fdim, inlet_facets)

    u_inlet = Function(V)

    def inlet_profile(x):
        # x has shape (3, N); rows are x, y, z coordinates
        # Centroid of inlet (approximate; refine if asymmetric)
        xc = x[0].mean()
        zc = x[2].mean()
        r2 = (x[0] - xc)**2 + (x[2] - zc)**2
        speed = U_max * np.maximum(0.0, 1.0 - r2 / R**2)
        values = np.zeros((3, x.shape[1]))
        values[1] = -speed  # y-component, negative = into domain
        return values

    u_inlet.interpolate(inlet_profile)
    bc_inlet = dirichletbc(u_inlet, inlet_dofs)
    return bc_inlet
\end{lstlisting}

\subsection*{Wall: no-slip}

\begin{lstlisting}
def make_wall_bc(V, mesh, tags):
    fdim = mesh.topology.dim - 1
    wall_facets = tags.find(1)
    wall_dofs = locate_dofs_topological(V, fdim, wall_facets)
    u_zero = Function(V)  # default initialises to zero
    return dirichletbc(u_zero, wall_dofs)
\end{lstlisting}

\subsection*{Outlet: zero-pressure Neumann}

The natural BC at the outlet is a zero-traction condition
$(-p\mathbf{I} + \mu\nabla\mathbf{u})\cdot\hat{\mathbf{n}} = \mathbf{0}$,
which is the default Neumann condition in the weak form. \emph{Do nothing}
at the outlet in the code; the weak form already encodes it.

However, since the pressure is only determined up to a constant, a
pressure nullspace must be handled explicitly (Section~\ref{sec:nullspace}).

\section{Weak form}

The symmetric weak form used by dolfinx (with the sign-flipped pressure):
\[
  a((\mathbf{u},p), (\mathbf{v},q))
    = \int_\Omega \nabla\mathbf{u}:\nabla\mathbf{v}\,dx
    - \int_\Omega p\,\nabla\cdot\mathbf{v}\,dx
    + \int_\Omega q\,\nabla\cdot\mathbf{u}\,dx
    = \int_\Omega \mu^{-1}\mathbf{f}\cdot\mathbf{v}\,dx
\]
where $\mathbf{f}=\mathbf{0}$ (no body force). The viscosity $\mu$ enters
as a scaling on the left-hand side when using dimensional coordinates.

\begin{lstlisting}
import ufl
from dolfinx.fem import Constant, form
from petsc4py import PETSc

def make_stokes_forms(mesh, V, Q, mu: float = 3.5e-3):
    """Build the blocked UFL forms for the Stokes problem.

    Parameters
    ----------
    mu : float
        Dynamic viscosity in Pa.s.

    Returns
    -------
    a : compiled blocked bilinear form (2x2 list)
    L : compiled blocked linear form (2-list)
    a_p11 : compiled pressure mass matrix for preconditioner
    """
    u, p = ufl.TrialFunction(V), ufl.TrialFunction(Q)
    v, q = ufl.TestFunction(V), ufl.TestFunction(Q)

    mu_c = Constant(mesh, PETSc.ScalarType(mu))
    f    = Constant(mesh, PETSc.ScalarType((0.0, 0.0, 0.0)))

    # Bilinear form (2x2 block structure)
    a_ufl = [
        [mu_c * ufl.inner(ufl.grad(u), ufl.grad(v)) * ufl.dx,
         ufl.inner(p, ufl.div(v)) * ufl.dx],
        [ufl.inner(ufl.div(u), q) * ufl.dx, None],
    ]
    # Linear form
    L_ufl = [ufl.inner(f, v) * ufl.dx, ufl.ZeroBaseForm((q,))]

    # Pressure mass matrix for block-diagonal preconditioner
    a_p11_ufl = ufl.inner(p, q) * ufl.dx

    return form(a_ufl), form(L_ufl), form(a_p11_ufl)
\end{lstlisting}

\textbf{Important.} The viscosity is included in the weak form as a
\texttt{Constant}, not absorbed into non-dimensionalisation, because the
CFD solve is in dimensional variables. Non-dimensionalisation is applied
later, at the synthetic-MRI and PINN stages.

\section{Pressure nullspace and solver}
\label{sec:solver}

The Stokes problem with a pure-Neumann pressure BC (our outlet) has
pressure defined only up to a constant. Without nullspace removal, MINRES
will stagnate or produce a drift-contaminated pressure field.

\textbf{Recommended solver: MINRES with block-diagonal preconditioner.}
This is the standard choice for symmetric indefinite saddle-point problems
of this size ($\sim 160\text{k}$ nodes $\times$ 4 unknowns $\approx 640\text{k}$ DOFs).

\begin{lstlisting}
from dolfinx.fem import (
    bcs_by_block, extract_function_spaces, locate_dofs_topological,
)
from dolfinx.fem.petsc import (
    LinearProblem, assemble_matrix, assemble_vector,
    apply_lifting, create_vector, set_bc,
)
from dolfinx.la.petsc import create_vector_wrap

def solve_stokes(mesh, V, Q, a, L, a_p11, bcs):
    """Solve the blocked Stokes system with MINRES + fieldsplit preconditioner.

    Handles the pressure nullspace explicitly.

    Returns
    -------
    u_h : dolfinx.fem.Function  velocity solution on V
    p_h : dolfinx.fem.Function  pressure solution on Q
    """
    a_p = [[a[0][0], None], [None, a_p11]]

    problem = LinearProblem(
        a,
        L,
        kind="nest",
        bcs=bcs,
        P=a_p,
        petsc_options_prefix="stokes_",
        petsc_options={
            "ksp_type":                   "minres",
            "ksp_rtol":                   1e-10,
            "ksp_max_it":                 500,
            "pc_type":                    "fieldsplit",
            "pc_fieldsplit_type":         "additive",
            "fieldsplit_0_ksp_type":      "preonly",
            "fieldsplit_0_pc_type":       "gamg",
            "fieldsplit_1_ksp_type":      "preonly",
            "fieldsplit_1_pc_type":       "jacobi",
        },
    )

    # Pressure nullspace: constant vector over the pressure subspace
    null_vec = create_vector(extract_function_spaces(L), "nest")
    null_vecs = null_vec.getNestSubVecs()
    null_vecs[0].set(0.0)   # velocity part: zero
    null_vecs[1].set(1.0)   # pressure part: constant
    null_vec.normalize()
    nsp = PETSc.NullSpace().create(vectors=[null_vec])
    problem.A.setNullSpace(nsp)

    # Mark preconditioner blocks as SPD for GAMG
    A00 = problem.A.getNestSubMatrix(0, 0)
    A00.setOption(PETSc.Mat.Option.SPD, True)
    P00 = problem.P_mat.getNestSubMatrix(0, 0)
    P00.setOption(PETSc.Mat.Option.SPD, True)
    P11 = problem.P_mat.getNestSubMatrix(1, 1)
    P11.setOption(PETSc.Mat.Option.SPD, True)

    u_h, p_h = problem.solve()

    converged = problem.solver.getConvergedReason()
    if converged <= 0:
        raise RuntimeError(
            f"Stokes solver did not converge. "
            f"PETSc reason code: {converged}. "
            f"Try reducing ksp_rtol or increasing ksp_max_it."
        )
    return u_h, p_h
\end{lstlisting}

\subsection*{Pressure sign correction after solve}

After solving, correct the pressure sign so it matches the physical
convention $p_{\text{physical}} = -p_{\text{dolfinx}}$:

\begin{lstlisting}
import numpy as np

def correct_pressure_sign(p_h) -> None:
    """Negate the dolfinx pressure field in-place to recover physical pressure."""
    p_h.x.array[:] = -p_h.x.array[:]
\end{lstlisting}

\subsection*{Pressure anchoring (alternative to nullspace)}

If iterative nullspace removal is insufficient (rare), anchor the pressure
at a single outlet degree of freedom to zero:

\begin{lstlisting}
from dolfinx.fem import locate_dofs_topological, dirichletbc

def make_pressure_anchor(Q, mesh, tags) -> "DirichletBC":
    """Pin pressure = 0 at one outlet DOF. Use as last resort if
    nullspace removal is insufficient."""
    outlet_facets = tags.find(3)
    p_dofs = locate_dofs_topological(Q, mesh.topology.dim - 1, outlet_facets)
    # Anchor only the first DOF to avoid over-constraining
    return dirichletbc(PETSc.ScalarType(0.0), p_dofs[:1], Q)
\end{lstlisting}

\section{WSS extraction}
\label{sec:wss}

WSS is the tangential traction on the wall:
\[
  \boldsymbol{\tau}_w
    = 2\mu\,\mathbf{D}(\mathbf{u})\cdot\hat{\mathbf{n}}
      - 2\mu\,\bigl(\hat{\mathbf{n}}^\top\mathbf{D}(\mathbf{u})\hat{\mathbf{n}}\bigr)\hat{\mathbf{n}},
  \qquad
  \mathbf{D} = \tfrac{1}{2}(\nabla\mathbf{u} + (\nabla\mathbf{u})^\top).
\]

In FEniCSx, WSS is computed as a boundary expression using UFL and
projected onto the wall facets.

\begin{lstlisting}
import ufl
from basix.ufl import element
from dolfinx import default_real_type
from dolfinx.fem import Expression, Function, functionspace
from dolfinx.fem.petsc import LinearProblem

def compute_wss(
    mesh, u_h, tags, mu: float = 3.5e-3
) -> "dolfinx.fem.Function":
    """Compute WSS magnitude field on the wall surface.

    Projects the WSS vector onto a DG-0 (piecewise constant) scalar
    space on wall facets, returning ||tau_w|| in Pa.

    The result is a Function on the full mesh DG-0 space; only the
    values at wall facets carry physical meaning.
    """
    fdim = mesh.topology.dim - 1
    n = ufl.FacetNormal(mesh)

    # Strain rate tensor
    D = 0.5 * (ufl.grad(u_h) + ufl.grad(u_h).T)

    # Traction vector on wall (using sign-corrected physical viscosity)
    traction = 2.0 * mu * ufl.dot(D, n)

    # Tangential component (subtract normal projection)
    traction_normal = ufl.dot(traction, n) * n
    wss_vector = traction - traction_normal

    # WSS magnitude
    wss_mag = ufl.sqrt(ufl.dot(wss_vector, wss_vector))

    # Project onto DG-0 space over wall facets
    # Use a DG-0 space (one value per cell, evaluated at wall)
    DG0_elem = element("DG", mesh.basix_cell(), degree=0,
                        dtype=default_real_type)
    S = functionspace(mesh, DG0_elem)
    wss_h = Function(S)

    # Use ds measure restricted to wall facets
    wall_facets = tags.find(1)
    ds_wall = ufl.Measure("ds", domain=mesh,
                           subdomain_data=tags, subdomain_id=1)

    v = ufl.TestFunction(S)
    u_trial = ufl.TrialFunction(S)

    # L2 projection of wss_mag onto S over wall facets
    # Note: lumped mass approximation (divide by facet area) via
    # explicit L2 projection
    a_proj = ufl.inner(u_trial, v) * ds_wall
    L_proj = ufl.inner(wss_mag, v) * ds_wall

    from dolfinx.fem import form as fem_form
    prob = LinearProblem(
        a_proj, L_proj,
        petsc_options={"ksp_type": "cg", "pc_type": "jacobi",
                       "ksp_rtol": 1e-12}
    )
    wss_h = prob.solve()
    return wss_h
\end{lstlisting}

\textbf{Alternative (simpler, lower accuracy): pointwise evaluation.}
For diagnostic purposes, evaluate the WSS expression at wall facet
midpoints using \texttt{Expression} and \texttt{eval}. This avoids the
projection solve but gives pointwise values rather than a spatially
consistent field. Use only for sanity checks.

\subsection*{WSS vector field (for direction-sensitive metrics)}

For OSI and vector WSS maps, export the full vector:

\begin{lstlisting}
def compute_wss_vector(mesh, u_h, tags, mu=3.5e-3):
    """Compute WSS vector (not just magnitude) on wall facets.

    Returns a Function on a DG-0 vector space (dim 3).
    """
    from basix.ufl import element
    from dolfinx import default_real_type
    from dolfinx.fem import functionspace, Function
    from dolfinx.fem import form as fem_form
    import ufl

    n = ufl.FacetNormal(mesh)
    D = 0.5 * (ufl.grad(u_h) + ufl.grad(u_h).T)
    traction = 2.0 * mu * ufl.dot(D, n)
    wss_vec = traction - ufl.dot(traction, n) * n

    DG0v = element("DG", mesh.basix_cell(), degree=0,
                   shape=(mesh.geometry.dim,), dtype=default_real_type)
    S3 = functionspace(mesh, DG0v)
    v  = ufl.TestFunction(S3)
    u_ = ufl.TrialFunction(S3)
    ds_wall = ufl.Measure("ds", domain=mesh,
                          subdomain_data=tags, subdomain_id=1)
    a_ = ufl.inner(u_, v) * ds_wall
    L_ = ufl.inner(wss_vec, v) * ds_wall
    from dolfinx.fem.petsc import LinearProblem
    prob = LinearProblem(
        a_, L_,
        petsc_options={"ksp_type": "cg", "pc_type": "jacobi",
                       "ksp_rtol": 1e-12}
    )
    return prob.solve()
\end{lstlisting}

\section{Output format: what the pipeline expects}

The downstream pipeline (synthetic MRI operator and PINN) reads CFD
output from HDF5 files via meshio or h5py. Export the following per case:

\begin{center}
\begin{tabular}{lll}
\toprule
\textbf{File} & \textbf{Contents} & \textbf{Format} \\
\midrule
\texttt{velocity.xdmf / .h5} & $\mathbf{u}$ on P2 nodes & XDMF + HDF5 \\
\texttt{pressure.xdmf / .h5} & $p$ on P1 nodes (sign-corrected) & XDMF + HDF5 \\
\texttt{wss\_mag.xdmf / .h5} & $\|\boldsymbol{\tau}_w\|$ on DG-0 cells & XDMF + HDF5 \\
\texttt{wss\_vec.xdmf / .h5} & $\boldsymbol{\tau}_w$ on DG-0 cells & XDMF + HDF5 \\
\texttt{mesh.xdmf / .h5} & Mesh geometry + boundary tags & XDMF + HDF5 \\
\texttt{metadata.json} & Physical params, Re, U, mesh path, solve time & JSON \\
\bottomrule
\end{tabular}
\end{center}

\begin{lstlisting}
import json
from dolfinx.io import XDMFFile
from mpi4py import MPI

def export_cfd_results(
    out_dir: str,
    mesh, tags, u_h, p_h, wss_mag_h, wss_vec_h,
    metadata: dict,
) -> None:
    """Write all CFD outputs to out_dir in XDMF+HDF5 format."""
    import pathlib
    out = pathlib.Path(out_dir)
    out.mkdir(parents=True, exist_ok=True)

    comm = MPI.COMM_WORLD

    for fname, func in [
        ("velocity",  u_h),
        ("pressure",  p_h),
        ("wss_mag",   wss_mag_h),
        ("wss_vec",   wss_vec_h),
    ]:
        with XDMFFile(comm, str(out / f"{fname}.xdmf"), "w") as f:
            f.write_mesh(mesh)
            func.x.scatter_forward()
            f.write_function(func)

    # Write mesh with boundary tags
    with XDMFFile(comm, str(out / "mesh.xdmf"), "w") as f:
        f.write_mesh(mesh)
        f.write_meshtags(tags, mesh.geometry)

    # Write metadata
    if comm.rank == 0:
        with open(out / "metadata.json", "w") as f:
            json.dump(metadata, f, indent=2)
\end{lstlisting}

\section{Full pipeline function}

\begin{lstlisting}
import time

def run_stokes_cfd(
    vtk_path: str,
    out_dir: str,
    mu: float = 3.5e-3,
    U_mean: float = 2e-5,
) -> None:
    """End-to-end Stokes CFD for one AnXplore case.

    Parameters
    ----------
    vtk_path : str
        Path to the AnXplore VTK mesh file (in mm).
    out_dir : str
        Directory to write all output files.
    mu : float
        Dynamic viscosity in Pa.s.
    U_mean : float
        Mean inlet velocity in m/s. Default gives Re ~ 0.10.
    """
    import tempfile, os
    from mpi4py import MPI

    t0 = time.time()

    # 1. Convert mesh
    with tempfile.NamedTemporaryFile(suffix=".xdmf", delete=False) as tmp:
        xdmf_tmp = tmp.name
    vtk_to_xdmf(vtk_path, xdmf_tmp)

    # 2. Load mesh
    mesh = load_mesh(xdmf_tmp)
    os.unlink(xdmf_tmp)
    os.unlink(xdmf_tmp.replace(".xdmf", ".h5"))

    # 3. Tag boundaries
    tags = tag_boundaries(mesh)

    # Sanity check
    assert len(tags.find(2)) > 0, "Inlet facets not found — check geometry"
    assert len(tags.find(3)) > 0, "Outlet facets not found — check geometry"

    # 4. Function spaces
    V, Q = make_function_spaces(mesh)

    # 5. Forms
    a, L, a_p11 = make_stokes_forms(mesh, V, Q, mu=mu)

    # 6. Boundary conditions
    bc_inlet = make_inlet_bc(V, mesh, tags, U_mean=U_mean)
    bc_wall  = make_wall_bc(V, mesh, tags)
    bcs = [bc_wall, bc_inlet]

    # 7. Solve
    u_h, p_h = solve_stokes(mesh, V, Q, a, L, a_p11, bcs)

    # 8. Correct pressure sign
    correct_pressure_sign(p_h)

    # 9. Compute WSS
    wss_mag_h = compute_wss(mesh, u_h, tags, mu=mu)
    wss_vec_h = compute_wss_vector(mesh, u_h, tags, mu=mu)

    t1 = time.time()

    # 10. Export
    rho = 1060.0  # kg/m^3
    L_char = 0.01628  # m
    Re = rho * U_mean * L_char / mu

    export_cfd_results(
        out_dir, mesh, tags, u_h, p_h, wss_mag_h, wss_vec_h,
        metadata={
            "vtk_source":  vtk_path,
            "mu_Pa_s":     mu,
            "rho_kg_m3":   rho,
            "U_mean_m_s":  U_mean,
            "Re":          Re,
            "L_char_m":    L_char,
            "wall_time_s": round(t1 - t0, 1),
            "dolfinx_version": "0.10",
        }
    )
    if MPI.COMM_WORLD.rank == 0:
        print(f"Done. Re={Re:.3f}, wall time={t1-t0:.1f}s -> {out_dir}")
\end{lstlisting}

\section{Known gotchas and failure modes}

\begin{enumerate}[noitemsep]

\item \textbf{\texttt{FunctionSpace} vs \texttt{functionspace}.}
  In dolfinx $\geq 0.7$, the constructor is the lowercase
  \texttt{functionspace(mesh, element)}. The CamelCase
  \texttt{FunctionSpace} import was removed. Using the old form raises
  an \texttt{ImportError} that is easy to miss.

\item \textbf{\texttt{basix.ufl.element} required.}
  From dolfinx~0.7 onward, elements must be created via
  \texttt{basix.ufl.element}, not \texttt{ufl.FiniteElement} or
  \texttt{ufl.VectorElement} (those are UFL legacy syntax that
  dolfinx~0.9+ no longer accepts). Always import from \texttt{basix.ufl}.

\item \textbf{Missing \texttt{create\_connectivity}.}
  Omitting \texttt{mesh.topology.create\_connectivity(fdim, tdim)}
  before \texttt{locate\_entities\_boundary} causes empty facet arrays
  with no error message. The boundary conditions are then silently
  not applied.

\item \textbf{Solver convergence failure with empty BCs.}
  If the inlet BC is not applied (empty inlet facet set), the system
  is singular and MINRES will not converge. The error message from
  PETSc (\texttt{DIVERGED\_BREAKDOWN} or similar) does not point to
  the BC. Always check \texttt{len(tags.find(2)) > 0} before solving.

\item \textbf{Pressure drift without nullspace.}
  Stokes with pure-Neumann pressure has a one-dimensional nullspace.
  If not removed, MINRES converges to a solution with arbitrary pressure
  offset. The velocity is correct (it is orthogonal to the nullspace),
  but the pressure and any derived quantity involving $p$ will be wrong.
  Always attach the nullspace or use a pressure anchor.

\item \textbf{Coordinate units.}
  AnXplore meshes are in millimetres. FEniCSx uses whatever units the
  coordinates are in. Failing to scale to metres means $\mu$ is
  dimensionally inconsistent with the geometry: viscous stresses will be
  $10^6\times$ too large or too small, and the solution will look
  physically unreasonable without an obvious error.

\item \textbf{WSS projection on zero-measure interior cells.}
  The DG-0 projection of WSS uses a \texttt{ds} measure restricted to
  wall facets (\texttt{subdomain\_id=1}). Cells not adjacent to wall
  facets get zero contribution, producing a zero WSS value rather than
  an error. Always mask interior cells when computing WSS statistics.

\item \textbf{Ghost update before export.}
  In parallel runs, call \texttt{func.x.scatter\_forward()} before
  writing to XDMF. In serial runs it is a no-op but should be included
  for correctness.

\item \textbf{XDMF sidecar \texttt{.h5} file.}
  Every \texttt{.xdmf} output also creates a \texttt{.h5} file. The
  meshio reader needs both; do not delete the \texttt{.h5}. When
  downstream code reads CFD output with meshio, pass the \texttt{.xdmf}
  path; meshio locates the \texttt{.h5} automatically.

\end{enumerate}

\section{Validation: Poiseuille analytical test}

Before running on AnXplore geometry, validate the solver on a straight
cylinder of radius $R=0.002\,\text{m}$, length $L=0.02\,\text{m}$,
where the analytical solution is:
\[
  u_y(r) = \frac{1}{4\mu}\frac{\Delta p}{L}(R^2 - r^2),
  \qquad
  \tau_w = \frac{R}{2}\frac{|\Delta p|}{L}.
\]
The PINN project uses $U_{\text{mean}} = 2\times10^{-5}\,\text{m/s}$, so
the target pressure drop across $L=0.02\,\text{m}$ is:
\[
  \Delta p = \frac{8\mu\bar{U}L}{R^2}
  = \frac{8 \times 3.5\times10^{-3} \times 2\times10^{-5} \times 0.02}{(0.002)^2}
  \approx 2.8\times10^{-4}\,\text{Pa}.
\]

Accept the solver if: (i) the centreline velocity matches $U_{\max} = 2\bar{U}$
to within 1\%; (ii) the WSS magnitude matches $\tau_w = 4\mu\bar{U}/R =
7\times10^{-5}\,\text{Pa}$ to within 2\%; (iii) $\|\nabla\cdot\mathbf{u}\|_{L^2}
< 10^{-10}$ (discrete incompressibility).

\section{Required imports summary}

\begin{lstlisting}
# All imports required for this module (dolfinx 0.10)
from mpi4py import MPI
from petsc4py import PETSc
import numpy as np
import ufl
from basix.ufl import element, mixed_element       # NOT ufl.FiniteElement
from dolfinx import default_real_type, la
from dolfinx.fem import (
    Constant, Function,
    bcs_by_block, dirichletbc,
    extract_function_spaces, form,
    functionspace,                                  # NOT FunctionSpace
    locate_dofs_topological,
)
from dolfinx.fem.petsc import (
    LinearProblem,
    apply_lifting, assemble_matrix, assemble_vector,
    create_vector, set_bc,
)
from dolfinx.io import XDMFFile
from dolfinx.la.petsc import create_vector_wrap
from dolfinx.mesh import (
    locate_entities_boundary, meshtags,
)
\end{lstlisting}

\end{document}
